\scanpage{68}
% archive sheet number: 67
\[
\rho_s=\sigma_1\sigma_2\qquad\rho_f=\sigma_0\sigma_1
\]
\[
\begin{aligned}
\rho_s^p=\mathrm{id}\ \inserted{\rho_s^d\ne\mathrm{id}\text{ si }0<d<p}
&\Longleftrightarrow\deleted{\unsure{\Phi}}F_p(\alpha)=0
&&F_p\text{ polynôme demi-cyclotomique}\\
\rho_f^q=\mathrm{id},\ \inserted{\rho_f^d\ne\mathrm{id}\text{ si }0<d<q}
&\Longleftrightarrow F_q(\beta)=0
&&F_q\quad\text{``}
\end{aligned}
\]
\uline{Forme quadratique \unsure{invariante|invariantes}} \deleted{\unsure{en} $x,y,z,t$}
\[
\underline{ax^2+by^2}+cz^2+\underline{uyz}+\underline{vzx}+\underline{wxy}+\underline{px+qy+rz}+d
\]
Conditions sur les coeff : Comme \unsure{les ctes|les cte} sont inv. OPS $d=0$

L'invariance p.r. à $\sigma_0,\sigma_1$ ds le plan des $xy$ $(z=0)$ donne $a=b=-p=-q$, $w=-\unsure{\beta|\beta^2}a$, d'autre part l'inv. par $\sigma_2$ i.e. symétrie en $y,z$ donne $c=b,v=w,r=q$ de sorte que l'on doit avoir
\[
f(x,y,z)=a\bigl[x^2+y^2+z^2\deleted{\unsure{-\beta xy-\beta xz}}\inserted{\unsure{-\beta x(y+z)|\beta x(y+z)}}-x-y-z\bigr]+uyz
\]
On va exprimer l'invariance p.r. à $\sigma_0\inserted{\text{ puis }}\sigma_1$, en notant qu'on sait déjà que la forme $uyz$ est inv. en $\sigma_0$, donc il faut \deleted{exprimer que} exprimer que
\[
\deleted{a[x^2+y^2-\unsure{\alpha|\beta}xy-x-y]}
\]
$a[\text{---}]$ est invariant, c'est le cas. Reste à \deleted{(\unsure{d'abord la restriction en y faisant} $z=0$) \unsure{est}} \deleted{\unsure{donc}} exprimer l'invariance p.r. à $\sigma_1$, on trouve \deleted{l'invariance p. $\sigma_0,\sigma_1$} la condition \deleted{$u=$}
\[
\begin{aligned}
&\deleted{u=-a(\beta+2(1+\alpha)+}\\
&u=-a(2+2(\alpha+\beta)+\alpha\beta)
\end{aligned}
\]
d'où la forme ou la \deleted{\unsure{fonction}} \inserted{\unsure{fon.|form.}} quadratique fondamentale
